\documentclass[11pt]{article}

\usepackage[margin=25mm]{geometry}

\usepackage{amsmath,amssymb,fontspec,unicode-math}

\IfFontExistsTF{DejaVu Serif}{\setmainfont{DejaVu Serif}}{\setmainfont{Latin Modern Roman}}

\IfFontExistsTF{Latin Modern Math}{\setmathfont{Latin Modern Math}}{\setmathfont{DejaVu Math TeX Gyre}}

\usepackage{xurl,hyperref}

\setlength{\parindent}{0pt}\setlength{\parskip}{6pt}

\begin{document}

\title{Learning Cocoercive Conservative Denoisers via Helmholtz Decomposition for Poisson Imaging Inverse Problems}

\author{ELVIA preliminary AI analysis}\date{21 October 2026}\maketitle

Preliminary assessment: Major.\par

Independent scientific review remains pending. Author adoption has not been requested.

The current paper retains a nonconvergent exact two-cycle inside its allowed ADMM parameter range, divergence under its general convex-fidelity assumptions, and an invalid identification of pointwise Jacobian splitting with a Helmholtz field decomposition. Its printed spectral-norm iteration also returns zero on a nonzero skew operator. These Major theoretical/derivation defects do not establish failure of the reported imaging experiments or of the independent proximal characterization under exact conservative/cocoercive assumptions.

\section*{Review scope}

Latest-version, claim-focused review of the cited author-paper passages and applicable current author-release paths. The stated counterexamples and scope controls were checked analytically; no author program, old verifier, training run or benchmark was executed.

\section*{Problem 1: The printed iteration does not compute spectral norm}

Source: Appendix A.5, p.28, Algorithm 1; p.29, Eqs.60,64.

Assessment: supported; impact: minor

\subsection*{Source claim (paraphrase)}

Algorithm 1 uses normalized \(Aq\) iteration and final \(q^\top Aq\) as a spectral-norm estimate.

\subsection*{Analysis}

For \(R=\begin{pmatrix}0&-1\\1&0\end{pmatrix}\) and \(A=2R\), \(\|A\|_2=2\), but \(q^\top Aq=0\) for every real \(q\). Since \(A\) is invertible, every normalization step is defined. The printed procedure returns zero regardless of iteration count. Its implementation in the paper's training was unavailable.

\subsection*{Finding and scope}

Algorithm 1 is mathematically false as a general spectral-norm estimator, including for the skew operator it is explicitly asked to measure. The actual training estimator and the provenance of published finite norm values remain unresolved; this is not a claim that every experimental value is wrong.

\subsection*{Proposed correction}

Use a singular-value iteration through \(A^\top A\), with nullspace and stopping conditions; a finite iteration count alone does not certify an arbitrary operator norm.

Both the current proceedings and latest arXiv v2 retain the Rayleigh-quotient iteration. The current release README and file tree provide test scripts, without identifying an operative spectral-training routine that could resolve the implementation question.

\section*{Problem 2: The allowed ADMM range admits an exact two-cycle}

Source: Theorem 4.1, p.7; Appendix A.6, pp.29–31, Eqs.67–76.

Assessment: supported; impact: major

\subsection*{Source claim (paraphrase)}

Theorem 4.1 claims full-sequence convergence for the printed gamma/t interval.

\subsection*{Analysis}

Take \(\beta=1,\gamma=0.1,t=0.2\), and let \(b_*\in(0,0.5)\) be the positive root of \(1.4b_*=0.8\tanh(2b_*)\), approximately \(0.33204524\). Define \[F_0(x)=0.1x^2-0.8\log\cosh x,\qquad F(x)=F_0(x)-b_*x,\qquad P=\operatorname{prox}_F,\qquad D=(P-0.8I)/0.2.\] Here \(F''\in[-0.6,0.2]\); thus \(F\) is smooth, coercive and real analytic, hence KL. Its proximal subproblem is strongly convex and \(P'\in[5/6,5/2]\), giving \(D'\in[1/6,8.5]\subset[0,10]\). The denoiser is conservative in one dimension and \(0.1\)-cocoercive, with \(D^t=P\). The paper’s constants \(r=L=9/14\) are valid bounds. Its characteristic polynomial at \(t=0.2\) is \(-0.0416<0\), so the parameter is allowed.

Use zero-count Poisson fidelity \(G(x)=2x+\iota_{[0,\infty)}(x)\), with \(\operatorname{prox}_G(z)=\max(z-2,0)\), and the stated zero initialization. If \(P_0=\operatorname{prox}_{F_0}\), the root identity and oddness give \(P_0(b_*)=2b_*\), \(P_0(-b_*)=-2b_*\), and \(P(z)=P_0(z+b_*)\). Therefore the printed ADMM equations produce \[u_k=0,\qquad(v_k,b_k)=\begin{cases}(2b_*,-2b_*),&k\text{ odd},\\(-2b_*,0),&k\text{ even}.\end{cases}\] The largest input to the fidelity prox is \(4b_*<2\), which verifies the asserted zero values inductively. The sequence does not converge and its consensus residual remains \(2b_*>0\). The claimed positive descent coefficient instead equals \(1/2-9/28-(9/14)^2=-23/98\).

\subsection*{Finding and scope}

The allowed universal parameter range in current Theorem 4.1 is disproved by an exact, well-defined Poisson instance satisfying its assumptions. This does not imply failure of the principal experiments at their different parameter choice. Restoring the descent coefficient alone does not supply the separate boundedness and joint-KL requirements.

\subsection*{Proposed correction}

A positive descent coefficient requires the additional restriction \(t<\min(t_0,\gamma/(1-\gamma))\). Boundedness and joint-KL conditions remain separate; experiments at another \(\gamma\) are not shown to cycle.

Current Theorem 4.1 and Appendix A.6 retain the same range and reversed comparison in the second Lipschitz branch; latest arXiv v2 has the same steps. The existing stricter sufficient range addresses the coefficient only, with full convergence requiring additional conditions.

\section*{Problem 3: General convex fidelity does not ensure bounded iterates}

Source: Theorems 4.1,4.4 and their convergence proofs.

Assessment: supported; impact: major

\subsection*{Source claim (paraphrase)}

Theorems \(4.1\) and \(4.4\) allow general lower-semicontinuous convex G with the listed F assumptions.

\subsection*{Analysis}

Let \(F(x)=0.1\sqrt{1+x^2}\), \(G(x)=-x\), \(\gamma=0.25,t=0.2\). The former is smooth, convex, coercive and analytic; the latter is proper closed convex as the general theorems allow. For each stated \(\beta\), put \(P=\operatorname{prox}_{F/\beta}\) and \(D=(P-0.8I)/0.2\), so \(D^t=P\). At \(\beta=1\), \(D'\in[6/11,1]\subset[0,4]\); at \(\beta=2\), \(D'\in[16/21,1]\subset[0,4]\). These maps are conservative and \(0.25\)-cocoercive. The ADMM parameter is within \(t<1/3\), and the PEGD parameter satisfies the printed step-size bound.

Since \(\operatorname{prox}_G(z)=z+1\), ADMM reduces to \(v_{k+1}=\operatorname{prox}_F(v_k+1)\). Proximal optimality and \(|F'|<0.1\) imply \(v_{k+1}-v_k>0.9\). For PEGD, the unit-parameter Moreau envelope is \({}^1G(x)=-x-1/2\); at \(\beta=2\), the increment similarly exceeds \(0.45\). Thus both sequences diverge. The respective objective derivatives are below \(-0.9\), so there is no finite stationary point to converge to. The affine fidelity is not the paper’s experimental Poisson fidelity: the contradiction concerns the literal general convex-fidelity quantifier.

\subsection*{Finding and scope}

The current general convergence statements do not follow from their listed assumptions. Coercivity of the denoiser potential alone cannot control a jointly unbounded-below objective. The result does not refute convergence for a separately justified Poisson objective or establish divergence of a learned network.

\subsection*{Proposed correction}

Require boundedness/coercivity of the relevant complete objective and valid joint-KL/step conditions. Safe numerical steps alone do not supply them.

The latest paper versions retain the general convex-fidelity wording and omitted boundedness/joint-KL conditions. Existing conditional sufficient revisions add assumptions; they are not a proof of the original unrestricted statement.

\section*{Problem 4: Pointwise Jacobian splitting is not a Helmholtz decomposition}

Source: Section 2.2, p.5, Eq.13; interpretation on p.6.

Assessment: supported; impact: major

\subsection*{Source claim (paraphrase)}

Section 2.2 identifies symmetric/skew Jacobian parts with Jacobians of conservative/divergence-free components.

\subsection*{Analysis}

For \(D(x,y)=(0,x^2)\), the symmetric Jacobian part is \(\begin{pmatrix}0&x\\x&0\end{pmatrix}\). If this were a field Jacobian, its first component would satisfy \(\partial_x d_1=0\) and \(\partial_y d_1=x\); their mixed derivatives would be \(0\) and \(1\), a contradiction. A smooth compact cutoff can preserve this local contradiction. Pointwise matrix splitting is therefore not a general integrable field decomposition.

\subsection*{Finding and scope}

The current pointwise matrix-to-field identification is false, even locally for a smooth field. The symmetric-Jacobian criterion for a conservative field on a suitable domain remains valid, and the independent proximal characterization is not refuted by this example.

\subsection*{Proposed correction}

Use the antisymmetric Jacobian as a local curl penalty. A genuine Helmholtz/Hodge field decomposition needs a domain, boundary conditions and a separate construction.

Current Section 2.2 and latest arXiv v2 retain the claimed Jacobian correspondence. Interpreting the skew penalty as a local nonconservativity measure avoids this false decomposition; it does not establish the stronger field or denoising interpretation.

\section*{Limits}

No author test/training program or old verifier was executed. The current release documents testing and does not identify the training spectral estimator needed to determine its experimental impact. The general convex-fidelity divergence example is not a Poisson imaging example. The two-cycle uses a permitted Poisson fidelity and zero initialization, but its parameter differs from the principal experimental setting.\par

\section*{Source identity}

NeurIPS 2025 official proceedings PDF (current publisher bytes, checked 2026-10-21)

\url{https://proceedings.neurips.cc/paper_files/paper/2025/file/c70bc769d6171ef8bb76e36e8d458205-Paper-Conference.pdf}

Reviewed PDF SHA-256: \nolinkurl{f0aec5a52afdf127a1b7d9c5427ebed751059b9b642d399205de1dc1dee25116}

This is the same research work as the predecessor, matched by title, author identity and official publisher channel. Its freshly retrieved exact PDF bytes and real retrieval time are recorded. Current target arXiv metadata and the relevant available author-version passages were checked separately; background works were excluded. Original report/source history is unchanged.

\end{document}
